%%% Pre aktiváciu nastavení odkomentovať riadok  (nie v tomto súbore,  ale
%%% v prekladanom zdrojovom súbore, umiestnenom tesne pred načítaním makra
%%% ZfSME):
%\input \ZFroot AKindle % nastavenia pre Amazon Kindle (č/b e-book A6)

%% https://groups.google.com/g/comp.text.tex/c/SbzyM6ZdVwo/m/HuDeL4vKPgsJ
%% (viď Anil Trivedi, 5. 10. 1993, 13:20:49)
\long\def\Ignore#1\endIgnore{\relax}
\def\endIgnore{\relax} % možné použiť priamo \Ignore ... \endIgnore
\def\IfIgnore{0}
\ifx\sdef\undeffined \input opmac \def\IfIgnore{1} \fi
\ifx\ZFroot\undeffined \def\ZFroot{} \fi

%% Platí posledný odkomentovaný riadok z tejto dvojice:
\sdef{pgs:a6k}{(90,120)mm} % skutočná veľkosť plochy e-ink (je menšia ako A6) 
\sdef{pgs:a6k}{(96,128)mm} % veľkosť plochy e-ink (o niečo väčšia ako skutočná)

\hyperlinks{\Grey}{\Grey}         % nastavenie šedých klikacích odkazov
\def\Capprefix{\localcolor\Grey}  % prvé veľké písmeno bude šedé 
\def\ObrPdf{1}              % č/b verzia vkladaných {obrázkov
\def\LogoZF{-1}             % čísla strán v strede päty (pre e-book)
\def\ZahlavieTlac{0}        % bez tlače záhlavia
\def\obrys#1{\vbox{\hrule\hbox{\vrule#1\vrule}\hrule}}
\def\TitSkip{-0} 
%\def\TitSkip{-.04} % odkomentovať len pre viac mien, alebo dlhý názov
%% PDF prehliadač v Kindle zobrazuje len jej viditeľnú časť, bez ohľadu na
%% Bouding Box, preto sa číslo strany síce vytlačí, ale bude biele
%% (neviditeľné len pre oko čiateľa); tak isto ako tenká vodorovná biela
%% čiara na začiatku strany:
\def\nopagenumbers{\footline={\localcolor\White\hss\tenrm\folio\hss}}
\def\KindleHR{{\localcolor\White\hrule}}
%% Platí posledný odkomentovaný riadok z tejto dvojice:
%\margins/1 a6  (3,3,3,11)mm % nastavenie formátu A6, okraje na 3mm,
\margins/1 a6k (3,3,3,11)mm % nastavenie formátu pre AKindle, okraje na 3mm,
                            % spodný 11mm (číslo strany v strede päty)

\ifnum\IfIgnore = 1 \Ignore \fi
\endinput \endIgnore

%%% Kalibrácia obrazovky Amazon Kindle A6 %%%

\shyph % zapína slovenské delenie slov vo formáte csplain/pdfcsplain
\sdef{pgs:a6k}{(90,120)mm} % skutočná veľkosť plochy e-ink (je menšia ako A6) 
\margins/1 a6k (3,3,3,11)mm % nastavenie formátu pre AKindle
\def\Titul{Kalibrácia obrazovky Amazon Kindle~A6}
\def\Autor{Marián Prochocký}
\input pdfuni % umožní vkladať PDF informácie v kódovaní UTF-16BE Unicode
\pdfunidef\tmp{\Autor}    \pdfinfo{/Author (\tmp)}
\pdfunidef\tmp{\Titul}    \pdfinfo{/Title (\tmp)}
\activettchar"
\input \ZFroot ZfSME2

\def\nopagenumbers{\footline={\localcolor\Grey\hss\tenrm\folio\hss}}
\def\KindleHR{{\localcolor\LightGrey\hrule}}


\TitulStrana

%% Nastavenie číslovania strán
\pageno=1 \PageNumbers %obnovenie číslovania strán od strany 1

\seccc \Titul

Nasledujúca strana slúži na kontrolu správneho nastavenia hodnôt pre rozmery
plochy s~e\=ink v~súbore "AKindle.tex". Pokiaľ sú hodnoty rozmerov nastavené
správne,  mali by byť  rozmery obdĺžnika nasledovné  (pričom  1\,pt $\doteq$
0.35146\,mm):

"\hsize" = \the\hsize\par
"\vsize" = \the\vsize\par

\noindent
čo by  mala byť  cca 84\,mm  šírka  obdĺžnika  (tomu  zodpovedá  "\hsize") a
106\,mm jeho výška (tomu zodpovedá "\vsize").

Pri tomto je  však  potrebné  zabezpečiť  zobrazovanie  správnych  proporcií
strán, ktorých vyplnenie sa začína medzerou (patitul a titul knihy)~-- tu sa
doplní neviditeľná (biela) vodorovná  čiara  na  začiatok  strany;  prípadne
chýba číslo strany~-- tu sa doplní jeho číslo neviditeľným (bielym)  písmom.
V~tomto dokumente je pre tieto prvky na titulnej  strane  zmenená  farba  na
šedú, takže sú viditeľné.
\NewPage

\obrys{\hskip\hsize \vbox to\vsize{}}
\NewPage

%% veľkosť plochy e-ink (o niečo väčšia ako skutočná)
\sdef{pgs:a6k}{(96,128)mm} 
\margins/1 a6k (3,3,3,11)mm % nastavenie formátu pre AKindle

V~súbore "AKindle.tex" je alternatívne aj  možnosť  nastaviť  o~niečo  väšie
rozmery strany, ako je fyzická veľkosť plochy  e\=ink.  Jedná  sa  o  rozmer
strany s~rovnakou šírkou strany, ako je použitá pre  formát  "eb"  a  "ebr",
t.~j. 96\,mm, čo zabezpečí aj vhodné nastavenie ostatných parametrov, ako aj
zlom odstavcov do riadkov s~rovnakou šírkou. Toto sa  používa  obyčajne  pri
zložitejšie formátovaných knihách. Pri tomto dôjde  automaticky  k~zmenšeniu
strany. Keďže rozmery sú len o~niečo väčšie,  aj  písmo  sa  zmenší  v~tomto
pomere ($90:96$, t.~j. \FR 15/16 , tak ako na tejto strane):

"\hsize" = \the\hsize\par
"\vsize" = \the\vsize\par

\noindent
Výsledok je na nasledujúcej strane,  kde má výsledný obdĺžnik  (zobrazený na
zariadení Amazon Kindle) rovnaké rozmey, ako obdĺžnik na strane 2.
\NewPage

\obrys{\hskip\hsize \vbox to\vsize{}}

\bye

